\documentclass[11pt, reqno]{article}

\usepackage{amsmath, amsthm, amssymb}
\usepackage{enumitem}
\usepackage{tcolorbox}
\usepackage{hyperref}
\usepackage{tikz}
\usepackage{tikz-cd}
\usepackage{pgfplots}
\pgfplotsset{compat=1.18}
\usetikzlibrary{arrows.meta}
\usepackage{mathrsfs}
\usepackage{fancyhdr}
\usepackage[bottom=0.75in, top=1in, left=0.5in, right=0.5in]{geometry}
\usepackage{array}   % for \newcolumntype macro
\newcolumntype{L}{>{$}l<{$}}

\theoremstyle{plain}
\newtheorem*{theorem}{Theorem}
\newtheorem*{proposition}{Proposition}
\newtheorem{exercise}{Exercise}
\newtheorem*{lemma}{Lemma}
\newtheorem*{corollary}{Corollary}

\theoremstyle{definition}
\newtheorem*{definition}{Definition}
\newtheorem*{example}{Example}

\theoremstyle{remark}
\newtheorem*{remark}{Remark}

\renewcommand{\phi}{\varphi}
\renewcommand{\epsilon}{\varepsilon}
\renewcommand{\emptyset}{\varnothing}

\newcommand{\RR}{\mathbb{R}}
\newcommand{\ZZ}{\mathbb{Z}}
\newcommand{\NN}{\mathbb{N}}
\newcommand{\CC}{\mathbb{C}}
\newcommand{\QQ}{\mathbb{Q}}

\DeclareMathOperator{\ima}{\text{im}}

\begin{document}

\topmargin=-40pt
\rhead{Henry Woodburn}
\lhead{Math 622}
\renewcommand{\headrulewidth}{1pt}
\renewcommand{\headsep}{20pt}
\thispagestyle{fancy}

{\Huge \bfseries \noindent Homework 7}

\begin{enumerate}
    \item[1.] Let $\omega_1, \dots \omega_p$, $\theta_1, \dots, \theta_p$ be 1-forms on a $d$-dimensional
    manifold $M$ such that 
    \[
        \sum_{i = 1}^p \omega_i \wedge \theta_i = 0.
    \]
    Since $\{\omega_i\}|_m$ is a set of $p$ linearly independent vectors for every $m \in M$, we can 
    complete it to a set $\{\omega_i\}_{i = 1}^d$ for a fixed $m \in M$ on an open set $U \subset M$
    containing $m$. 

    Then we can express each $\theta_i$ as
    \[
        \theta_i = \sum_{j = 1}^d A_{i,j}\omega_j,
    \]
    $j = 1, \dots, p$. To show the coefficients are smooth, choose vector fields $V_1, \dots, V_d$ on $U$ such 
    that $\omega_j(V_i) = \delta_{j}^i$. Then 
    \[
        A_{i,j} = \theta_i(V_j),
    \]
    so $A_{i,j}$ is smooth since by definition, any one form applied to a smooth vector field yields a 
    smooth function. 

    Now we write
    \begin{align*}
        &0 = \sum_{i = 1}^p \theta_i \wedge \omega_i = \sum_{i = 1}^p\left(\sum_{j = 1}^d A_{j,i}\omega_j\right)
        \wedge \omega_i\\
        &= \sum_{i=1}^p \sum_{j = 1}^d A_{j,i}\omega_j \wedge \omega_i = \sum_{i,j = 1}^p A_{j,i} \omega_j \wedge \omega_i
        + \sum_{i = 1}^p \sum_{j = p + 1}^d A_{j,i} \omega_j \wedge \omega_i\\
        &= \sum_{i < k} A_{j,i}\omega_j \wedge \omega_i + \sum_{i = 1}^p \sum_{j = p + 1}^d A_{j,i} \omega_j \wedge \omega_i\\
        &= \sum_{i \leq k}(A_{i,j} - A_{j,i})\omega_i \wedge \omega_j + \sum_{i = 1}^p \sum_{j = p + 1}^d -A_{j,i} \omega_i \wedge \omega_j.
    \end{align*}

    Since $\omega_i \wedge \omega_j$, $i < j$, forms a basis for the space of $2$-forms, and $i < j$ for every
    element of the last sum, we have both 
    \[
        A_{j,i} = 0\ \text{for} j = p+ 1, \dots, d, i = 1, \dots, p
    \]
    and
    \[
        A_{j,i} = A_{i,j}\ \text{for} 1 \leq i,j \leq p.
    \]
    Then in particular, we can write
    \[
        \theta_j = \sum_{i = 1}^p A_{i,j} \omega_i
    \]
    where $A_{i,j}$ are smooth.

    \item[2.] Let $X$ be a smooth vector field, $m$ a covariant $s$-tensor field. Let $p \in M$. 
    Choose local coordinates $\frac{\partial}{\partial x_i}$, $1 = 1, \dots, n$ such 
    that $X = \frac{\partial}{\partial x_1}$. Then by definition, we have
    \[
        \mathcal{L}_X \omega|_p = \frac{d}{dt}|_{t = 0}(\theta_t^* \omega)|_p,
    \]
    where $\theta_t$ is the flow of $X$. Then $\mathcal{L}_X$ acts on a tensor
    field by taking the partial derivate with respect to $x_1$. 
    Let $\{u^i\}$ be a dual frame to $\{x_i\}$. For $s$ vector
    field $Y_j$, $j = 1, \dots, s$, we first have
    \begin{align*}
        \omega(Y_1, \dots, Y_s) = \omega_{i_1, \dots, i_s}u^{i_1}\otimes \cdots\otimes u^{i_s}
        \left(\left(Y^{1, j_1}x_{j_1}\right)\otimes \cdots\otimes \left(Y^{s, j_s}x_{j_s}\right)\right)\\
        = \omega_{i_1, \dots, i_s}u^{i_1, \dots, i_s}\left(Y^{1, j_1} \cdots Y_{s, j_s} x_{j_1, \dots, j_s}\right)\\
        = \omega_{i_1, \dots, i_s}Y_{1, i_1}\cdots Y_{s, i_s}.
    \end{align*}

    Then we apply the product rule, so that 
    \begin{align*}
        \mathcal{L}_X(\omega(Y_1, \dots, Y_s)) = \left(\frac{\partial}{\partial x_1}\omega_{i_1, \dots, i_s}\right)Y_{1, i_1}\cdots Y_{s, i_s}
        + \omega_{i_1, \dots, i_s}\left(\frac{\partial}{\partial x_1} Y_{1, i_1}\right) Y_{2, i_2}\cdots Y_{s , i_s}\\
        = \mathcal{L}_X(\omega)(Y_1, \dots, Y_s) + A(\mathcal{L}_X Y_1, \dots, Y_s) + \cdots + A(Y_1, \dots, \mathcal{L}_X Y_s).
    \end{align*}

    To finish, we use that $\mathcal{L}_X f = X(f)$ and $\mathcal{L}_X(Y) = [X,Y]$ and move terms so that
    \[
        \mathcal{L}_X\omega(Y_1, \dots, Y_s) = X(\omega(Y_1, \dots, Y_s)) - 
    \omega([X, Y_1], \dots, Y_s) - \cdots - \omega(Y_1, \dots, [X, Y_s])
    \]

<<<<<<< HEAD
    \item[3.] We will make use of the formula 
    \[
        i_{X}(\omega^1 \wedge \cdots \wedge \omega^k) = \sum{i = 1}^k (-1)^{i - 1} \omega^i(X)\wedge \cdots
        \wedge \hat{\omega^i}\wedge \cdots \wedge \omega^k
    \]
    which is proven in Lee. 

    Here it suffices to suppose $\omega$ and $\nu$ are decomposable. 
=======
    \item[3.] We will use the formula from Lee which states that 
    \begin{equation}\label{wedge}
        i_X(\omega^1 \wedge \omega^k) = \sum_{i = 1}^k (-1)^{i-1}\omega^i(X)\omega^i\wedge 
        \cdots \wedge \hat{\omega^i} \wedge \cdots \wedge \omega^k.
    \end{equation}
    To prove this formula, we let $X = X_0$, and let $X_1, \dots, X_{k - 1}$ be $k-1$ vector fields.
    We first know 
    \[
        i_X(\omega^1 \wedge \cdots \wedge \omega^k)(X_1,\dots, X_{k-1})
        = (\omega^1\wedge \cdots\wedge \omega^k)(X_0, \dots, X_{k-1}) = \det(\omega^j(X_{i})).
    \]
    Next we notice that the right side of $\ref{wedge}$ is exactly the expansion of this determinate 
    by the first column. This proves the formula.

    Since every $s$-form is a linear combination of decomposable $s$-forms, by bilinearity of the 
    wedge product, it suffices to prove the lemma for decomposable forms. Let $omega = \omega^1 \wedge \cdots \wedge \omega^s$
    and $\nu = \omega^{s+1}\wedge \cdots \wedge \omega^{s + t}$ be $s$ and $t$ forms. Then
    \begin{align*}
        &i_X{\omega \wedge \nu} = i_X(\wedge^1 \wedge \cdots \wedge \omega^{s + t})\\
        &= \sum_{i = 1}^{s+t}(-1)^{i-1}\omega^i(X)\omega^1 \wedge \cdots \wedge \hat{\omega^i}\wedge \cdots \wedge \omega^{s + t}\\
        &= \sum_{i = 1}^s (-1)^{i-1}\omega^i(X)\omega^1 \wedge \cdots \wedge \hat{\omega^i}\wedge \cdots \wedge \omega^{s + t}
        + \sum_{i = s+1}^{s + t} (-1)^{i-1}\omega^i(X)\omega^1 \wedge \cdots \wedge \hat{\omega^i}\wedge \cdots \wedge \omega^{s + t}\\
        &= \left(\sum_{i = 1}^s (-1)^{i-1}(-1)^{i-1}\omega^i(X)\omega^1 \wedge \cdots \wedge \hat{\omega^i}\wedge \cdots \wedge \omega^{s}\right)
        \wedge \left(\omega^{s+1}\wedge \cdots\wedge\omega^{s+t}\right)\\
        & + \sum_{i = 1}^{t}(-1)^{s+i-1}\left(\omega^1\wedge \cdots\wedge\omega^s\right)
        \wedge\left(\omega^{s+1}\wedge\cdots\wedge\hat{\omega^i}\wedge\cdots\wedge\omega^{s+t}\right)\\
        &= (i_X \omega)\wedge \nu + (-1)^s\omega\wedge(i_X\nu)
    \end{align*}

    and we are done. 

    \item[4.] First assume $\omega$ is a $p$-form on $M\times N$ such that both conditions 
    $i_X \omega = 0$ and $\mathcal{L}_X\omega = 0$ for all $X$ with $d\pi(X) = 0$. We let 
    $x_1,\dots, x_s$ and $x_{s+1},\dots, x_{s+t}$ be local coordinates for $M$ and $N$ in some neighborhood, and let the corresponding
    dual bases have upper indices. The last condition tells us that
    \[
        X = \sum_{i = 1}^s X_i(m,n)\frac{\partial}{\partial x_i}.
    \]

    Suppose first that $\omega = x^I$ for 
    some $I$ a multi-index of size $p$ by linearity. Choosing $X = \frac{\partial}{\partial x_i}$, we can write 
    \[
        0 = i_X\omega = \sum_{j = 1}^{s+t} (-1)^{j-1}m^{I_j}(\frac{\partial}{\partial x_i})m^{I_1}\wedge\cdots\wedge
        \hat{m^{I_j}}\wedge\cdots\wedge m^{I_{s+t}}
        = (-1)^{j-1}m^i(\frac{\partial}{\partial x_i}).
    \]
    But the last term would be nonzero if any $I_j \leq s$. Then $\omega$ must be a sum of 
    wedges of forms only from the cotangent space of $N$. 

    Now we show the coefficients of $\omega$ must also not depend on $M$. 

    First note that for an arbitrary vector field $V$, we have
    \[
        d\pi[X,V] = [d\pi X, d\pi V] = 0
    \]
    if $d\pi(X) = 0$. Then by a previous exercise,
    \[
        0 = \mathcal{L}_X\omega(V_1, \dots, V_p) = X(\omega(V_1,\dots, V_p)) - \omega([X,V_1], \dots, V_p)
        - \cdots - \omega(V_1, \dots, [X, V_p]) = X(\omega(V_1, \dots, V_p)),
    \]

    using the result mentioned above for lie brackets, since this means $i_{[X,V]}\omega = 0$.

    Now with 
    \[
        \omega = \sum_{I} a_I(m,n) m^I,
    \]

    we choose vector fields $V_1, \dots, V_p$ so that $\omega(V_1, \dots, V_p) = a_I$ for any given
    index $I$. Then we have
    \[
        0 = X(\omega(V_1, \dots, V_p)) = X(a_I(m,n))
    \]
    for any $X$ as above. This means the $a_I$ do not depend on inputs in $M$. Then 
    we can just let $\alpha = \sum_I a_I m^I$, with all $m^{I_i}$ in $T^*N$, so that 
    \[
        \omega(V_1, \dots, V_p) = \omega(\{0\}\times d\pi(V_1),\dots,\{0\}\times d\pi(V_p))
        = \alpha(d\pi(V_1),\dots,d\pi(V_p)) = \pi^*\alpha(V_1,\dots,V_p).
    \]
    
    Conversely, suppose $\omega = \pi^*\alpha$ for $\alpha$ a $p$-form on $N$. Then 
    $\omega = \sum_{I} a_I(n) dm^I$, where the $dm^I$ all come from $N$, where we view
    the $a_I(n)$ as functions on $M\times N$ which are constant across $M$, and where the elements
    of $T^*N$ are naturally considered as elements of $T*(M\times N)$. Then it is obvious
    that if $d\pi(X) = 0$, $i_X(\omega) = 0$, since $X$ only takes derivatives in $M$. Similar
    to the opposite direction, we can use the formula for the Lie derivative to deduce that
    $\mathcal{L}_X\omega = X\circ\omega$, and that this equals zero since the coefficients of
    $\omega$ only depend on values in $N$. 

    \item[5.] (a.) We apply Cartan's magic formula:
    \[
        \mathcal{L}_X \Omega = i_X(d\Omega) + d(i_X\Omega).
    \]

    Since $\Omega$ is a top-form, $d\Omega = 0$. Then $d(i_X \Omega) = \mathcal{L}_X \Omega = \text{div}_\Omega X \Omega$
    by definition.

    (b.) Let $M = \RR^n$, $X = \sum_{i} X^i \frac{\partial}{\partial x^i}$, $\Omega = dx^1 \wedge \cdots \wedge dx^n$.
    Then 
    \begin{align*}
        \text{div}_\Omega X \Omega = d(i_X\Omega) = d\left(\sum_{i = 1}^n (-1)^{i-1} X^i
        dx^1 \wedge \cdots \wedge \hat{dx^i}\wedge\cdots\wedge dx^n\right)\\
        = sum_{j = 1}^n \sum_{i = 1}^n \frac{\partial X^j}{\partial x^i}(-1)^{i-1} dx_j\wedge dx^1 \wedge \cdots \wedge \hat{dx^i}\wedge\cdots\wedge dx^n\\
        = \sum_{i=1}^n \frac{\partial x_i}{\partial x^i}dx^1\wedge\cdots \wedge dx^n
        = \left(\sum_{i = 1}^n \frac{\partial X_i}{\partial x^i}\right) \Omega.
    \end{align*}

    Then $\text{div}_\Omega X = \sum \frac{\partial X_i}{\partial x^i}$. 
>>>>>>> 20d82943e97e28887e835582274edf57856fc70e

\end{enumerate}

\end{document}